% Motor: XeLaTeX, fijado en .latexmkrc y en la receta del editor.
\documentclass[aspectratio=169,10pt,c]{beamer}
% Sin atajos de babel: así las comillas " del código se ven tal cual.
\PassOptionsToPackage{es-noshorthands}{babel}

% !TEX root = slides.tex
% Preambulo comun: estilo UTP, colores, fuentes, plantillas de beamer y
% comandos de apoyo didactico. Se carga desde slides.tex tras \documentclass.

% La opción "c" de beamer reparte el espacio libre 1:1.5 (arriba:abajo), por lo
% que el contenido queda por encima del centro. Igualamos ambos rellenos.
\makeatletter
\define@key{beamerframe}{c}[true]{%
  \beamer@frametopskip=0pt plus 1fill\relax%
  \beamer@framebottomskip=0pt plus 1fill\relax%
  \beamer@frametopskipautobreak=0pt plus .5\paperheight\relax%
  \beamer@framebottomskipautobreak=0pt plus .5\paperheight\relax%
}
\makeatother

\usepackage{graphicx}
\usepackage{tikz}
\usepackage{xcolor}
\usepackage{iftex}

\ifPDFTeX
  \usepackage[utf8]{inputenc}
  \usepackage[T1]{fontenc}
  \usepackage{lmodern}          % tipografia escalable para el respaldo con pdflatex
  \renewcommand{\familydefault}{\sfdefault}
\else
  \usepackage{fontspec}
  \IfFontExistsTF{Arial}
    {\setsansfont{Arial}[BoldFont={Arial Bold}]}
    {\setsansfont{Helvetica}[BoldFont={Helvetica Bold}]}
\fi

\usepackage[spanish,es-nodecimaldot]{babel}
\usepackage{microtype}
\usepackage{booktabs}
\usepackage{amssymb}
\usepackage{array}
\usetikzlibrary{arrows.meta,positioning,fit,calc,shapes.geometric}

\graphicspath{{assets/}}

% ---------- Identidad visual UTP ----------
\definecolor{utpRojo}{HTML}{C42B35}
\definecolor{utpNegro}{HTML}{111111}
\definecolor{utpGris}{HTML}{5F5F5F}
\definecolor{utpGrisMedio}{HTML}{9A9A9A}
\definecolor{utpGrisClaro}{HTML}{F1F1F1}
\definecolor{utpNaranja}{HTML}{C1670B}

\usetheme{default}
\setbeamercolor{normal text}{fg=utpNegro,bg=white}
\setbeamercolor{structure}{fg=utpRojo}
\setbeamercolor{frametitle}{fg=utpNegro,bg=white}
\setbeamercolor{title}{fg=utpNegro}
\setbeamercolor{subtitle}{fg=utpNegro}
\setbeamercolor{author}{fg=utpNegro}
\setbeamercolor{date}{fg=utpNegro}
\setbeamercolor{block title}{fg=white,bg=utpNegro}
\setbeamercolor{block body}{fg=utpNegro,bg=utpGrisClaro}
\setbeamercolor{alertblock title}{fg=white,bg=utpRojo}
\setbeamercolor{alertblock body}{fg=utpNegro,bg=utpRojo!7}
\setbeamercolor{exampleblock title}{fg=white,bg=utpGris}
\setbeamercolor{exampleblock body}{fg=utpNegro,bg=utpGris!8}

\setbeamerfont{title}{series=\bfseries,size=\fontsize{19.5}{22.5}\selectfont}
\setbeamerfont{subtitle}{size=\fontsize{11.5}{13.5}\selectfont}
\setbeamerfont{date}{series=\bfseries,size=\fontsize{11.5}{13.5}\selectfont}
\setbeamerfont{author}{size=\fontsize{10.5}{12.5}\selectfont}
\setbeamerfont{frametitle}{series=\bfseries,size=\large}
\setbeamerfont{block title}{series=\bfseries,size=\small}
\setbeamerfont{block body}{size=\normalsize}

\setbeamertemplate{navigation symbols}{}
\setbeamertemplate{headline}{}
\setbeamertemplate{blocks}[default]
\setbeamertemplate{itemize item}{\color{utpRojo}\small\raise.15ex\hbox{$\blacktriangleright$}}
\setbeamertemplate{itemize subitem}{\color{utpGris}\scriptsize$\bullet$}
\setbeamertemplate{enumerate item}{\color{utpRojo}\insertenumlabel.}
\setbeamertemplate{background canvas}{\color{white}\rule{\paperwidth}{\paperheight}}

% Título de diapositiva: negrita negra a la izquierda, logo UTP a la derecha
\setbeamertemplate{frametitle}{%
  \begin{tikzpicture}[remember picture,overlay]
    \node[anchor=north east,inner sep=0pt]
      at ([xshift=-7mm,yshift=-5.5mm]current page.north east)
      {\includegraphics[width=17mm]{utp-logo.png}};
  \end{tikzpicture}%
  \vspace{1.2mm}%
  \begin{minipage}[t]{.80\textwidth}\raggedright\insertframetitle\end{minipage}\par
  \vspace{1.1mm}{\color{utpRojo}\rule{\textwidth}{1.1pt}}
}

% Solo numeración discreta; la plantilla UTP no lleva pie de página
\setbeamertemplate{footline}{%
  \hfill{\color{utpGrisMedio}\scriptsize\insertframenumber}\hspace{6mm}\vspace{2.5mm}%
}

% Portada UTP
\setbeamertemplate{title page}{%
  \begin{tikzpicture}[remember picture,overlay]
    \node[anchor=center,inner sep=0pt] at (current page.center) {%
      \includegraphics[width=\paperwidth,height=\paperheight]{utp-cover-bg.jpeg}%
    };
    \node[anchor=north,align=center,text=utpNegro,inner sep=0pt,
      text width=145mm,font=\usebeamerfont{title}]
      at ([yshift=-25.2mm]current page.north) {\inserttitle};
    \node[anchor=north,align=center,text=utpNegro,inner sep=0pt,
      text width=142mm,font=\usebeamerfont{subtitle}]
      at ([yshift=-38.2mm]current page.north) {\insertsubtitle};
    \node[anchor=north,align=center,text=utpNegro,inner sep=0pt,
      font=\usebeamerfont{date}]
      at ([yshift=-50.4mm]current page.north) {\insertdate};
    \node[anchor=north,align=center,text=utpNegro,inner sep=0pt,
      font=\usebeamerfont{author}]
      at ([yshift=-58.4mm]current page.north) {\insertauthor};
  \end{tikzpicture}%
}

% ---------- Comandos de apoyo didáctico ----------
\newcommand{\concepto}[2]{%
  \begin{block}{#1}\large #2\end{block}}
\newcommand{\pregunta}[1]{%
  \begin{alertblock}{Decide y explica}#1\end{alertblock}}
\newcommand{\analogia}[1]{%
  \begin{exampleblock}{En palabras simples}\small #1\end{exampleblock}}
\newcommand{\enlarealidad}[1]{%
  \begin{block}{Cómo se ve en la realidad}\small #1\end{block}}
\newcommand{\ojo}[1]{%
  \begin{alertblock}{Ojo con esto}\small #1\end{alertblock}}
\newcommand{\apoyo}{{\normalfont\footnotesize\color{utpGris}~· material de apoyo}}
\newcommand{\marca}[1]{\tikz[baseline=(m.base)]{%
  \node[rounded corners=1.5pt,fill=utpGrisClaro,draw=utpGrisMedio,inner xsep=3pt,inner ysep=1.5pt,
        font=\footnotesize] (m) {#1};}}

\tikzset{
  caja/.style={draw=utpNegro, very thick, rounded corners=1.5pt, fill=white,
    align=center, minimum height=0.72cm, inner sep=5pt, font=\small},
  cajaazul/.style={caja, draw=utpGris, fill=white},
  cajaverde/.style={caja, draw=utpNegro, fill=utpGrisClaro},
  cajanaranja/.style={caja, draw=utpRojo!75, fill=utpRojo!7},
  flecha/.style={-{Stealth[length=2.3mm]}, very thick, draw=utpNegro},
  flechavuelta/.style={{Stealth[length=2.3mm]}-{Stealth[length=2.3mm]}, very thick, draw=utpNegro}
}

% ---------- Notación pata de gallo (crow's foot) ----------
% La rotación indica hacia dónde está la entidad: 0 derecha, 180 izquierda,
% 90 arriba, -90 abajo. La pata se abre hacia la entidad del lado "muchos".
\tikzset{
  relacion/.style={thick, draw=utpNegro},
  pics/uno/.style={code={%
    \draw[thick,utpNegro,line cap=round] (0,-1.5mm) -- (0,1.5mm);}},
  pics/muchos/.style={code={%
    \draw[thick,utpNegro,line cap=round,line join=round]
      (2.6mm,1.5mm) -- (0,0) -- (2.6mm,-1.5mm);
    \draw[thick,utpNegro,line cap=round] (0,0) -- (2.6mm,0);}},
}
\newcommand{\simbuno}{\tikz[baseline=-0.6ex]{\draw[relacion](0,0)--(7mm,0);\pic at (4.4mm,0){uno};}}
\newcommand{\simbmuchos}{\tikz[baseline=-0.6ex]{\draw[relacion](0,0)--(7mm,0);\pic at (4.4mm,0){muchos};}}

\usepackage{colortbl}

\title{Transacciones}
\subtitle{\texttt{START TRANSACTION}, \texttt{COMMIT} y \texttt{ROLLBACK} en MySQL y SQL Server}
\date{Semana 13 - Sesión 26}
\author{Rel Guzman}

% Nombre de tabla, siempre en naranja
\newcommand{\tbl}[1]{\textcolor{utpNaranja}{\texttt{#1}}}
% Palabra reservada de SQL dentro de un bloque de código
\newcommand{\kw}[1]{\textcolor{utpRojo}{\textbf{#1}}}
% Comentario dentro de un bloque de código
\newcommand{\com}[1]{\textcolor{utpGrisMedio}{#1}}
% Texto de apoyo en gris
\newcommand{\gris}[1]{\textcolor{utpGris}{#1}}
% Insignia de práctica
\newcommand{\ins}[1]{\tikz[baseline=-0.5ex]{\node[rounded corners=2pt,fill=utpRojo,text=white,
  font=\bfseries\scriptsize,inner xsep=5pt,inner ysep=2.5pt] {#1};}}
% Etiqueta oscura: nombre de un operador o de un gestor
\newcommand{\etiqueta}[1]{\tikz[baseline=-0.5ex]{\node[rounded corners=2pt,fill=utpNegro,text=white,
  font=\bfseries\scriptsize,inner xsep=6pt,inner ysep=3pt] {#1};}}
% Número de paso
\newcommand{\npieza}[1]{\tikz[baseline=-0.6ex]{\node[circle,fill=utpRojo,text=white,
  font=\bfseries\scriptsize,inner sep=1.6pt,minimum size=4.2mm] {#1};}}
% Celda o texto resaltado en rojo claro
\newcommand{\mal}[1]{{\setlength{\fboxsep}{1.2pt}\colorbox{utpRojo!18}{#1}}}
% Lo importante, resaltado en amarillo dentro de un bloque de código
\newcommand{\dentro}[1]{{\setlength{\fboxsep}{1.5pt}\colorbox{yellow!45}{#1}}}
% Fila de una tabla resaltada
\newcommand{\fila}{\rowcolor{yellow!45}}
\newcommand{\filagris}{\rowcolor{utpGrisClaro}}
% Marcas de sí y no
\newcommand{\si}{$\checkmark$}
\newcommand{\no}{\textcolor{utpRojo}{$\times$}}
% Bloques de código sobre fondo gris: normal, chico y mínimo
\newsavebox{\sqlbuf}
\newenvironment{sqlcode}[1][.86\textwidth]
  {\begin{center}\begin{lrbox}{\sqlbuf}\begin{minipage}{#1}\ttfamily\footnotesize\vspace{6pt}\raggedright}
  {\vspace{6pt}\end{minipage}\end{lrbox}\colorbox{utpGrisClaro}{\usebox{\sqlbuf}}\end{center}}
\newenvironment{sqlchico}[1][.86\textwidth]
  {\begin{center}\begin{lrbox}{\sqlbuf}\begin{minipage}{#1}\ttfamily\scriptsize\vspace{5pt}\raggedright}
  {\vspace{5pt}\end{minipage}\end{lrbox}\colorbox{utpGrisClaro}{\usebox{\sqlbuf}}\end{center}}
\newenvironment{sqlmini}[1][.86\textwidth]
  {\begin{center}\begin{lrbox}{\sqlbuf}\begin{minipage}{#1}\ttfamily\fontsize{7}{8.3}\selectfont\vspace{4pt}\raggedright}
  {\vspace{4pt}\end{minipage}\end{lrbox}\colorbox{utpGrisClaro}{\usebox{\sqlbuf}}\end{center}}
\newenvironment{sqlmicro}[1][.86\textwidth]
  {\begin{center}\begin{lrbox}{\sqlbuf}\begin{minipage}{#1}\ttfamily\fontsize{6.4}{7.5}\selectfont\vspace{4pt}\raggedright}
  {\vspace{4pt}\end{minipage}\end{lrbox}\colorbox{utpGrisClaro}{\usebox{\sqlbuf}}\end{center}}
% Ancho de una columna menos el borde del bloque de código
\newcommand{\col}{\dimexpr\textwidth-2\fboxsep\relax}
% Sangría dentro de un bloque de código
\newcommand{\ind}{\hspace*{1.2em}}
% Mensaje de error dentro de un bloque de código
\newcommand{\err}[1]{\textcolor{utpRojo}{#1}}
% Anatomía de una instrucción: cada pieza del código con su pregunta al lado
\tikzset{
  pieza/.style={fill=utpGrisClaro,font=\ttfamily\footnotesize,anchor=west,inner sep=5pt,
    minimum height=0.62cm},
  para/.style={font=\small,anchor=west,text=utpNegro}
}
\newcommand{\piezax}{7.6}
\newcommand{\pieza}[3]{%
  \node[pieza] at (0,#1) {#2};
  \node[para] at (\piezax,#1) {\textcolor{utpRojo}{$\blacktriangleleft$}\ #3};}
% Diapositiva de sección: una palabra bien grande
\newcommand{\seccion}[2][1]{%
\begin{frame}[plain]
  \begin{tikzpicture}[remember picture,overlay]
    \node[anchor=north east,inner sep=0pt]
      at ([xshift=-10mm,yshift=-8mm]current page.north east)
      {\includegraphics[width=43mm]{utp-logo.png}};
    \node[anchor=center,inner sep=0pt] at (current page.center)
      {\resizebox{100mm}{!}{\bfseries\color{utpNegro}#2}};
    \ifnum#1=1
    \node[anchor=north,inner sep=0pt] at ([yshift=-11mm]current page.center)
      {\color{utpRojo}\rule{100mm}{1.6pt}};
    \fi
  \end{tikzpicture}
\end{frame}}
% Cabecera de columna dentro de una tabla en monoespaciado
\newcommand{\cab}[1]{\normalfont\bfseries #1}
% Resultado de una consulta: \res{columnas}{cabecera}{filas}
\newcommand{\res}[3]{%
  \centering{\scriptsize\color{utpGris}resultado}\\[3pt]
  \footnotesize\ttfamily
  \begin{tabular}{#1}
  \toprule
  #2 \\
  \midrule
  #3
  \bottomrule
  \end{tabular}}
% Lo mismo, cuando debajo sigue más contenido del frame
\newcommand{\resg}[3]{{\res{#1}{#2}{#3}\par}}
% Tabla pequeña con su nombre arriba: \mini{nombre}{columnas}{cabecera}{filas}
\newcommand{\mini}[4]{%
  \begin{tabular}{@{}c@{}}
  \tbl{#1}\\[3pt]
  {\scriptsize\ttfamily
  \begin{tabular}{#2}
  \toprule
  #3 \\
  \midrule
  #4
  \bottomrule
  \end{tabular}}
  \end{tabular}}
% Una tabla con título gris, alineada arriba: \tabla{título}{columnas}{cabecera}{filas}
\newcommand{\tabla}[4]{%
  \begin{tabular}[t]{@{}c@{}}
  {\scriptsize\color{utpGris}#1}\\[3pt]
  {\footnotesize\ttfamily
  \begin{tabular}[t]{#2}
  \toprule
  #3 \\
  \midrule
  #4
  \bottomrule
  \end{tabular}}
  \end{tabular}}
% La pregunta del frame, en una caja naranja centrada
\newcommand{\preg}[1]{{\centering\tikz\node[caja,cajanaranja,font=\normalsize,inner sep=6pt]{#1};\par}\vspace{4pt}}
% Forma de lo que devuelve una consulta: celdas con borde
\newcommand{\forma}[2]{{\ttfamily\footnotesize\renewcommand{\arraystretch}{1.25}%
  \begin{tabular}{#1}\hline #2\end{tabular}}}
% Estilos de los diagramas de esta sesión
\tikzset{
  codigo/.style={font=\ttfamily\footnotesize,inner sep=3pt},
  sub/.style={codigo,draw=utpRojo,very thick,fill=yellow!45,rounded corners=2pt},
  decision/.style={diamond,aspect=2.6,draw=utpRojo,very thick,fill=utpRojo!7,
    font=\scriptsize,align=center,inner sep=1.5pt},
  paso/.style={caja,font=\scriptsize,minimum height=.6cm,inner sep=4pt}
}
\begin{document}
{\setbeamertemplate{footline}{}\maketitle}
\begin{frame}[c]
\begin{block}{Logro de la sesión}
Al finalizar, podrás registrar una venta con una \textbf{transacción}, que se guarda \textbf{entera o nada} con \texttt{COMMIT} y \texttt{ROLLBACK}, en MySQL y SQL Server.
\end{block}
\end{frame}
% =======================================================================
% Revisión de la sesión anterior
% ---------------------------------------------------------------------
\begin{frame}{Revisión: ¿qué hace \texttt{OVER}?}
\renewcommand{\arraystretch}{0.85}
\small Los 3 pagos de hoy y su promedio, de dos formas:
\vspace{0.2cm}
\begin{columns}[T,onlytextwidth]
\column{.42\textwidth}
\centering\etiqueta{Sin OVER}
\begin{sqlchico}[\col]
\kw{SELECT AVG}(monto) \kw{AS} promedio\\
\kw{FROM} pago \kw{WHERE} id >= 6;
\end{sqlchico}
\resg{r}{\cab{promedio}}{35.000000 \\}
\vspace{6pt}
\small Junta las filas: queda \textbf{1}.
\column{.54\textwidth}
\centering\etiqueta{Con OVER}
\begin{sqlchico}[\col]
\kw{SELECT} id, monto, \dentro{\kw{AVG}(monto) \kw{OVER} ()} \kw{AS} promedio\\
\kw{FROM} pago \kw{WHERE} id >= 6 \kw{ORDER BY} id;
\end{sqlchico}
\resg{rrr}{\cab{id} & \cab{monto} & \cab{promedio}}{6 & 15.00 & \cellcolor{yellow!45}35.000000 \\ 7 & 35.00 & \cellcolor{yellow!45}35.000000 \\ 8 & 55.00 & \cellcolor{yellow!45}35.000000 \\}
\vspace{6pt}
\small No junta: quedan las \textbf{3}, cada una con el promedio al lado.
\end{columns}
\vspace{0.3cm}
\concepto{\texttt{OVER}}{\small Calcula sobre varias filas \textbf{sin juntarlas}.}
\end{frame}
% ---------------------------------------------------------------------
\begin{frame}{Revisión: lo que va dentro de \texttt{OVER}}
\renewcommand{\arraystretch}{0.85}
\begin{center}
\small
\renewcommand{\arraystretch}{1.3}
\begin{tabular}{@{}ll@{}}
\toprule
\bfseries Se escribe & \bfseries Calcula sobre \\
\midrule
\texttt{OVER ()} & todas las filas \\
\texttt{OVER (PARTITION BY metodo)} & cada grupo por separado \\
\texttt{OVER (ORDER BY id)} & las filas en orden, sumando lo anterior \\
\bottomrule
\end{tabular}
\end{center}
\vspace{0.1cm}
\begin{columns}[c,onlytextwidth]
\column{.55\textwidth}
\begin{sqlchico}[\col]
\kw{SELECT} id, monto,\\
\ind\dentro{\kw{SUM}(monto) \kw{OVER} (\kw{ORDER BY} id)} \kw{AS} acumulado\\
\kw{FROM} pago \kw{WHERE} id >= 6 \kw{ORDER BY} id;
\end{sqlchico}
\column{.42\textwidth}
\centering
\tabla{resultado}{rrr}{\cab{id} & \cab{monto} & \cab{acumulado}}{6 & 15.00 & 15.00 \\ 7 & 35.00 & 50.00 \\ 8 & 55.00 & 105.00 \\}\\[4pt]
{\footnotesize\gris{15, luego 15 + 35, luego 50 + 55}}
\end{columns}
\vspace{0.2cm}
{\centering\small También van con \texttt{OVER}: \texttt{ROW\_NUMBER}, \texttt{RANK}, \texttt{DENSE\_RANK}, \texttt{LAG} y \texttt{LEAD}.\par}
\end{frame}
% =======================================================================
\seccion{Transacciones}
\begin{frame}{Caso de estudio: venta de entradas para una feria}
\begin{columns}[T,onlytextwidth]
\column{.54\textwidth}
\small La feria \textbf{Fogón 2026} dura tres días (4, 5 y 6 de diciembre) y vende su entrada en línea.
\vspace{0.05cm}
\begin{itemize}\setlength\itemsep{2pt}\footnotesize
  \item Hay \textbf{una sola entrada}, válida para los tres días: S/ 45, con \textbf{500 disponibles}.
  \item El \textbf{visitante} compra en la página principal: indica su \textbf{DNI}, su \textbf{nombre} y su \textbf{correo}, y paga con tarjeta.
  \item Cada DNI compra \textbf{una sola entrada}.
  \item Cada compra \textbf{descuenta una} de las entradas disponibles.
  \item Si el pago no se aprueba, \textbf{no} se registra la compra \textbf{ni} se descuenta la entrada.
  \item El \textbf{administrador} fija el precio y las entradas disponibles, abre o cierra la venta y consulta las ventas.
\end{itemize}
\column{.42\textwidth}
\centering\renewcommand{\arraystretch}{0.9}
\begin{tikzpicture}
  \tikzset{ent/.style={draw=utpNegro,very thick,fill=white,inner sep=0pt,anchor=north}}
  \node[ent] (e) at (0,0) {\begin{tabular}{@{}l@{}}
    \multicolumn{1}{@{}c@{}}{\colorbox{utpNegro}{\makebox[2.9cm]{\color{white}\bfseries\footnotesize entrada}}}\\[1pt]
    {\ttfamily\scriptsize\ \ nombre}\\
    {\ttfamily\scriptsize\ \ precio}\\
    {\ttfamily\scriptsize\ \ \colorbox{yellow!45}{disponibles}}\\
    \\[-6pt]
  \end{tabular}};
  \node[ent] (v) at (0,-2.55) {\begin{tabular}{@{}l@{}}
    \multicolumn{1}{@{}c@{}}{\colorbox{utpNegro}{\makebox[2.9cm]{\color{white}\bfseries\footnotesize venta}}}\\[1pt]
    {\ttfamily\scriptsize\ \ dni \gris{\normalfont(UNIQUE)}}\\
    {\ttfamily\scriptsize\ \ nombre}\\
    {\ttfamily\scriptsize\ \ correo}\\
    {\ttfamily\scriptsize\ \ precio}\\
    {\ttfamily\scriptsize\ \ referencia\_pago}\\[2pt]
  \end{tabular}};
  \draw[very thick,utpNegro] (e.south) -- (v.north);
  \node[font=\scriptsize\bfseries,anchor=west] at ([xshift=2pt,yshift=-6pt]e.south) {1};
  \node[font=\scriptsize\bfseries,anchor=west] at ([xshift=2pt,yshift=6pt]v.north) {muchas};
\end{tikzpicture}
\end{columns}
\vspace{0.2cm}
\concepto{En esta sesión}{\small Una compra cambia \textbf{las dos tablas}: resta 1 en \tbl{entrada} y agrega una fila en \tbl{venta}.}
\end{frame}
% ---------------------------------------------------------------------
\begin{frame}{Vender una entrada: dos cambios}
\small Quedan \textbf{500 entradas disponibles}. Luis Rojas compra la suya:
\vspace{0.3cm}
\begin{center}
\begin{tikzpicture}
  \tikzset{et/.style={caja,minimum width=3.4cm,minimum height=1.5cm,text width=3.2cm}}
  \node[et] (a) at (0,0) {\textbf{Portada}\\[3pt]\footnotesize Luis escribe su DNI, su nombre y su correo y paga \textbf{S/ 45}};
  \node[et] (b) at (4.6,0) {\textbf{Pasarela de pago}\\[3pt]\footnotesize\textcolor{utpRojo}{\checkmark} aprobado\\{\ttfamily PAGO-0001}};
  \node[et,cajanaranja] (c) at (9.2,0) {\textbf{Base de datos}\\[3pt]\footnotesize registra la venta:\\\textbf{dos cambios}};
  \draw[flecha] (a)--(b); \draw[flecha] (b)--(c);
\end{tikzpicture}
\end{center}
\vspace{0.15cm}
\begin{columns}[T,onlytextwidth]
\column{.3\textwidth}
\centering\npieza{1}\ \small Descontar \textbf{1 entrada disponible}\\[5pt]
\mini{entrada}{r}{\cab{disponibles}}{500 $\rightarrow$ \colorbox{yellow!45}{499} \\}
\column{.68\textwidth}
\centering\npieza{2}\ \small Guardar la \textbf{venta}\\[5pt]
\mini{venta}{llll}{\cab{dni} & \cab{nombre} & \cab{correo} & \cab{referencia\_pago}}{\fila 70123456 & Luis Rojas & luis.rojas@example.com & PAGO-0001 \\}
\end{columns}
\vspace{0.2cm}
\concepto{Idea clave}{\small Los dos cambios van \textbf{juntos}: o se hacen los dos, o ninguno.}
\end{frame}
% ---------------------------------------------------------------------
\begin{frame}{¿Y si el segundo cambio falla?}
\small Luis \textbf{ya compró} su entrada: quedan 499. Intenta comprar otra con el mismo DNI. Cada persona compra solo una: \texttt{dni} es \texttt{UNIQUE}.
\vspace{0.2cm}
\begin{sqlcode}[.9\textwidth]
\kw{UPDATE} entrada \kw{SET} disponibles = disponibles - 1 \kw{WHERE} id = 1;\hfill\textcolor{utpRojo}{\checkmark}\\[3pt]
\kw{INSERT INTO} venta (dni, nombre, correo, precio, referencia\_pago)\\
\kw{VALUES} (\mal{'70123456'}, 'Luis Rojas', 'luis.rojas@example.com',\\
\ind 45.00, 'PAGO-0002');\hfill\no\\[3pt]
\scriptsize\err{-- ERROR 1062 (23000): Duplicate entry '70123456' for key 'venta.dni'}
\end{sqlcode}
\vspace{0.1cm}
\begin{center}
\begin{tikzpicture}
  \node[paso,text width=3.6cm] (a) at (0,0) {\npieza{1} Se descontó 1 entrada\\{\ttfamily disponibles = \mal{498}}};
  \node[paso,text width=3.6cm,draw=utpRojo] (b) at (5,0) {\npieza{2} La venta \textbf{no} se guardó};
  \node[caja,cajanaranja,font=\small,text width=3.4cm] (c) at (10,0) {Se perdió una entrada que \textbf{nadie compró}};
  \draw[flecha] (a)--(b); \draw[flecha] (b)--(c);
\end{tikzpicture}
\end{center}
\vspace{0.2cm}
\ojo{Cada instrucción se guarda \textbf{apenas termina}. El \texttt{UPDATE} ya quedó guardado cuando falló el \texttt{INSERT}.}
\end{frame}
% ---------------------------------------------------------------------
\begin{frame}{Transacción: todo o nada}
\concepto{Transacción}{\small Un grupo de cambios que se guardan \textbf{todos juntos} o \textbf{ninguno}.}
\vspace{0.35cm}
\begin{center}
\begin{tikzpicture}
  \node[paso,fill=utpNegro,text=white,font=\scriptsize\bfseries\ttfamily] (ini) at (0,0) {START TRANSACTION};
  \node[paso,text width=2.6cm,align=left,font=\scriptsize\ttfamily] (cam) at (3.9,0) {\npieza{1} \kw{UPDATE} \ldots\\[2pt]\npieza{2} \kw{INSERT} \ldots};
  \node[decision] (d) at (7.4,0) {¿todo salió\\bien?};
  \node[caja,cajanaranja,font=\scriptsize,text width=2.7cm] (c) at (10.7,0.85) {\texttt{\textbf{COMMIT}}\\se guardan \textbf{los dos}};
  \node[caja,font=\scriptsize,text width=2.7cm] (r) at (10.7,-0.85) {\texttt{\textbf{ROLLBACK}}\\se deshacen \textbf{los dos}};
  \draw[flecha] (ini)--(cam);
  \draw[flecha] (cam)--(d);
  \draw[flecha] (d.north) |- node[above,pos=.75,font=\scriptsize\bfseries]{sí} (c.west);
  \draw[flecha] (d.south) |- node[below,pos=.75,font=\scriptsize\bfseries]{no} (r.west);
\end{tikzpicture}
\end{center}
\vspace{0.35cm}
\begin{center}
\footnotesize
\renewcommand{\arraystretch}{1.2}
\begin{tabular}{@{}lll@{}}
\toprule
 & \bfseries MySQL & \bfseries SQL Server \\
\midrule
Empezar & \texttt{START TRANSACTION;} & \texttt{\dentro{BEGIN} TRANSACTION;} \\
Guardar todo & \texttt{COMMIT;} & \texttt{COMMIT;} \\
Deshacer todo & \texttt{ROLLBACK;} & \texttt{ROLLBACK;} \\
\bottomrule
\end{tabular}
\end{center}
\end{frame}
% ---------------------------------------------------------------------
\begin{frame}{La venta a Luis: \texttt{COMMIT}}
\renewcommand{\arraystretch}{0.85}
\begin{sqlcode}[.8\textwidth]
\dentro{\kw{START TRANSACTION}};\\
\kw{UPDATE} entrada \kw{SET} disponibles = disponibles - 1 \kw{WHERE} id = 1;\\
\kw{INSERT INTO} venta (dni, nombre, correo, precio, referencia\_pago)\\
\kw{VALUES} ('70123456', 'Luis Rojas', 'luis.rojas@example.com',\\
\ind 45.00, 'PAGO-0001');\\
\dentro{\kw{COMMIT}};
\end{sqlcode}
\vspace{0.1cm}
\begin{columns}[T,onlytextwidth]
\column{.34\textwidth}
\centering
{\scriptsize\ttfamily\kw{SELECT} disponibles \kw{FROM} entrada;}\\[6pt]
\tabla{resultado}{r}{\cab{disponibles}}{\cellcolor{yellow!45}499 \\}
\column{.62\textwidth}
\centering
{\scriptsize\ttfamily\kw{SELECT} dni, nombre, referencia\_pago \kw{FROM} venta \kw{ORDER BY} id;}\\[6pt]
\tabla{resultado}{lll}{\cab{dni} & \cab{nombre} & \cab{referencia\_pago}}{\fila 70123456 & Luis Rojas & PAGO-0001 \\}
\end{columns}
\vspace{0.3cm}
\concepto{\texttt{COMMIT}}{\small Todo salió bien: los \textbf{dos cambios} quedan guardados a la vez.}
\end{frame}
% ---------------------------------------------------------------------
\begin{frame}{Luis otra vez: \texttt{ROLLBACK}}
\renewcommand{\arraystretch}{0.85}
\begin{columns}[c,onlytextwidth]
\column{.6\textwidth}
\begin{sqlchico}[\col]
\dentro{\kw{START TRANSACTION}};\\
\kw{UPDATE} entrada \kw{SET} disponibles = disponibles - 1 \kw{WHERE} id = 1;\\
\kw{SELECT} disponibles \kw{FROM} entrada;\hfill\npieza{1}\\[4pt]
\kw{INSERT INTO} venta (dni, nombre, correo, precio,\\
\ind referencia\_pago)\\
\kw{VALUES} (\mal{'70123456'}, 'Luis Rojas', 'luis.rojas@example.com',\\
\ind 45.00, 'PAGO-0002');\\
\scriptsize\err{-- ERROR 1062: Duplicate entry '70123456' for key 'venta.dni'}\\[4pt]
\scriptsize\dentro{\kw{ROLLBACK}};\\
\kw{SELECT} disponibles \kw{FROM} entrada;\hfill\npieza{2}
\end{sqlchico}
\column{.37\textwidth}
\centering
\begin{tikzpicture}
  \node (a) at (0,0) {\npieza{1}\ \tabla{pendiente}{r}{\cab{disponibles}}{\mal{498} \\}};
  \node (b) at (0,-2.3) {\npieza{2}\ \tabla{después del ROLLBACK}{r}{\cab{disponibles}}{\cellcolor{yellow!45}499 \\}};
  \draw[flecha] (a)--(b);
\end{tikzpicture}
\end{columns}
\vspace{0.3cm}
\concepto{\texttt{ROLLBACK}}{\small Algo falló: se deshace \textbf{todo}. La entrada vuelve a estar disponible y no queda ninguna venta a medias.}
\end{frame}
% ---------------------------------------------------------------------
\begin{frame}{\ins{Práctica 1} Una venta con transacción}
\small En cada gestor ejecuta \textbf{tres scripts, en orden}. Están en \texttt{ejemplos/practicas/mysql/} y \texttt{ejemplos/practicas/sqlserver/}:
\begin{center}
\footnotesize
\renewcommand{\arraystretch}{1.25}
\begin{tabular}{@{}cll@{}}
\toprule
\bfseries Paso & \bfseries Script & \bfseries Qué hace \\
\midrule
\npieza{1} & \texttt{practica\_1\_modelo\_fisico.sql} & crea la base \texttt{feria} y sus tablas \\
\npieza{2} & \texttt{practica\_1\_registros.sql} & carga la entrada: 500 disponibles y ninguna venta \\
\npieza{3} & \texttt{practica\_1\_transaccion.sql} & la venta a Luis, en \textbf{dos partes} \\
\bottomrule
\end{tabular}
\end{center}
\vspace{0.05cm}
\small El paso \npieza{3} se ejecuta \textbf{por partes}: selecciona una parte, ejecútala y mira el resultado.
\vspace{0.1cm}
\begin{center}
\footnotesize
\renewcommand{\arraystretch}{1.2}
\begin{tabular}{@{}cll@{}}
\toprule
\bfseries Parte & \bfseries Qué hace & \bfseries \texttt{disponibles} \\
\midrule
1 & La venta a Luis, con \texttt{COMMIT} & 500 $\rightarrow$ 499 \\
2 & Luis otra vez, con \texttt{ROLLBACK} & 499 $\rightarrow$ 498 (pendiente) $\rightarrow$ 499 \\
\bottomrule
\end{tabular}
\end{center}
\vspace{0.05cm}
\ojo{En la parte 2, el \texttt{INSERT} de Luis falla a propósito. Workbench se detiene ahí: selecciona las líneas que siguen y ejecútalas. Para repetir la práctica, vuelve a ejecutar los pasos \npieza{1} y \npieza{2}.}
\end{frame}
% ---------------------------------------------------------------------
\begin{frame}{Transacciones en MySQL, SQL Server y SQLite}
\begin{center}
\footnotesize
\renewcommand{\arraystretch}{1.2}
\begin{tabular}{@{}>{\raggedright\bfseries}p{3.4cm}>{\raggedright}p{3.3cm}>{\raggedright}p{3.0cm}>{\raggedright\arraybackslash}p{3.0cm}@{}}
\toprule
 & \bfseries MySQL & \bfseries SQL Server & \bfseries SQLite \\
\midrule
Empezar & \texttt{START TRANSACTION;} & \texttt{BEGIN TRANSACTION;} & \texttt{BEGIN TRANSACTION;} \\
Guardar todo & \texttt{COMMIT;} & \texttt{COMMIT;} & \texttt{COMMIT;} \\
Deshacer todo & \texttt{ROLLBACK;} & \texttt{ROLLBACK;} & \texttt{ROLLBACK;} \\
\midrule
Sin transacción & \multicolumn{3}{l}{cada instrucción se guarda \textbf{apenas termina}} \\
¿Un error la deshace? & \multicolumn{3}{l}{\textbf{no}: tú decides \texttt{ROLLBACK}} \\
La venta a Luis & \multicolumn{3}{l}{mismo resultado: 500 $\rightarrow$ 499 $\rightarrow$ 498 $\rightarrow$ 499} \\
\midrule
\texttt{CREATE TABLE} dentro & \mal{se guarda al instante}: \texttt{ROLLBACK} no la quita & \texttt{ROLLBACK} la quita & \texttt{ROLLBACK} la quita \\
\bottomrule
\end{tabular}
\end{center}
\vspace{0.05cm}
\concepto{¿Funcionan igual?}{\small \textbf{Sí} con \texttt{INSERT}, \texttt{UPDATE} y \texttt{DELETE}: todo o nada. Solo cambia cómo se empieza, y en MySQL no se crean tablas dentro de una transacción.}
\end{frame}
% ---------------------------------------------------------------------
\begin{frame}[plain]
  \begin{tikzpicture}[remember picture,overlay]
    \node[anchor=north west,font=\bfseries\fontsize{24}{28}\selectfont]
      at ([xshift=12mm,yshift=-9mm]current page.north west) {Preguntas};
    \node[anchor=north east,inner sep=0pt]
      at ([xshift=-10mm,yshift=-8mm]current page.north east)
      {\includegraphics[width=43mm]{utp-logo.png}};
    \node[anchor=center,inner sep=0pt]
      at ([yshift=-5mm]current page.center)
      {\includegraphics[width=78mm]{preguntas-ilustracion.png}};
  \end{tikzpicture}
\end{frame}
\begin{frame}{Cierre}
\begin{center}
\small
\renewcommand{\arraystretch}{1.12}
\begin{tabular}{>{\raggedright\bfseries}p{3.6cm}>{\raggedright\arraybackslash}p{9.4cm}}
\toprule
Transacción & Varias instrucciones que se guardan \textbf{todas o ninguna}. \\
Empezar & MySQL: \texttt{START TRANSACTION}; SQL Server y SQLite: \texttt{BEGIN TRANSACTION}. \\
\texttt{COMMIT} & Confirma y guarda los cambios. \\
\texttt{ROLLBACK} & Deshace los cambios: todo vuelve a como estaba. \\
Sin transacción & Cada instrucción se guarda apenas termina: si la siguiente falla, no hay cómo deshacerla. \\
\bottomrule
\end{tabular}
\end{center}
\end{frame}

\begin{frame}{Referencias}
\footnotesize
\begin{itemize}\setlength\itemsep{7pt}
\item Silberschatz, A., Korth, H. F., \& Sudarshan, S. (2014). \textit{Fundamentos de bases de datos} (6.ª ed.). McGraw-Hill.
\item Oracle. MySQL Reference Manual: \href{https://dev.mysql.com/doc/refman/8.0/en/commit.html}{\textit{START TRANSACTION, COMMIT, and ROLLBACK Statements}}.
\item Microsoft. Transact-SQL: \href{https://learn.microsoft.com/sql/t-sql/language-elements/begin-transaction-transact-sql}{\textit{BEGIN TRANSACTION}} y \href{https://learn.microsoft.com/sql/t-sql/language-elements/rollback-transaction-transact-sql}{\textit{ROLLBACK TRANSACTION}}.
\item SQLite. \href{https://www.sqlite.org/lang_transaction.html}{\textit{Transaction}}.
\end{itemize}
\end{frame}
% Cierre UTP
{
  \setbeamertemplate{background canvas}{%
    \includegraphics[width=\paperwidth,height=\paperheight]{gracias-final.pdf}%
  }
  \setbeamertemplate{footline}{}
  \begin{frame}[plain]
  \end{frame}
}

\end{document}
